\chapter{Definiciones y abreviaciones del Negocio}
\label{ap:definiciones-negocio}

Este apéndice reúne los conceptos que se emplean en F.A.M.A. para describir
la gestión de grabaciones, el entrenamiento de modelos y la revisión de
predicciones. Las definiciones se presentan en el contexto de la
bioacústica y del diagnóstico acústico de motores.

\section{Definiciones del dominio}

\begin{description}
    \item[Administrador] Usuario que gestiona cuentas y asignaciones de
    recursos y puede designar el modelo predeterminado. La operación de la
    infraestructura corresponde al responsable técnico del nodo.

    \item[Arquitectura híbrida] Organización que combina el procesamiento
    local de audios y modelos con servicios de almacenamiento y persistencia
    en la nube.

    \item[Bioacústica] Estudio de los sonidos producidos por los seres vivos.
    En el proyecto se aplica al reconocimiento de especies de aves a partir
    de grabaciones.

    \item[Checkpoint] Archivo que conserva los pesos de un modelo en un
    estado determinado y puede incluir su configuración y procedencia.
    Permite recuperar el modelo con un predictor compatible.

    \item[Clase acústica] Categoría que se desea reconocer, como una especie
    de ave o una condición de funcionamiento de un motor. La etiqueta es
    el valor que identifica esa categoría en un registro.

    \item[Clasificación acústica] Tarea de asignar una grabación a una de
    las clases contempladas para un dominio acústico.

    \item[Confianza de predicción] Puntuación que el modelo asigna a la clase
    predicha. Su interpretación como probabilidad de acierto requiere una
    evaluación de calibración.

    \item[Conjunto de datos (\textit{dataset})] Colección organizada de
    grabaciones, etiquetas y metadatos destinada al desarrollo o evaluación
    de modelos. El acceso puede ser privado o concederse a otros usuarios
    mediante las asignaciones correspondientes.

    \item[Curación de datos] Revisión de audios, etiquetas y procedencia
    para decidir qué aportes son aptos para incorporarse a un conjunto de
    datos.

    \item[Diagnóstico acústico de motores] Análisis de las señales sonoras
    del motor para clasificar grabaciones según las condiciones de operación
    o falla contempladas por el modelo.

    \item[Ensamble] Modelo compuesto que combina las salidas de varios
    modelos para obtener una predicción conjunta.

    \item[Entrenamiento] Proceso en el que se ajustan los pesos de un modelo
    utilizando ejemplos de la partición de entrenamiento.

    \item[Espectrograma Mel] Representación numérica de la energía de una
    grabación a lo largo del tiempo y en bandas de frecuencia en escala Mel.
    Puede visualizarse como imagen y utilizarse como entrada del modelo.

    \item[Evaluación] Medición del desempeño de un modelo mediante la
    comparación entre sus predicciones y las etiquetas de referencia de
    una partición identificada.

    \item[Exactitud (\textit{accuracy})] Proporción de ejemplos evaluados
    cuya clase predicha coincide con la etiqueta de referencia.

    \item[Exhaustividad (\textit{recall})] Proporción de ejemplos de una
    clase que el modelo identifica correctamente como pertenecientes a ella.

    \item[Fuga de datos] Uso directo o indirecto de información reservada
    para evaluar el modelo durante su entrenamiento o ajuste. Puede ocurrir
    si fragmentos de una misma grabación quedan en particiones distintas.

    \item[Grupo de origen] Conjunto de grabaciones relacionadas, por ejemplo
    por provenir del mismo audio original o de un grabador identificado.
    Se conserva dentro de una misma partición para reducir la fuga de datos.

    \item[Inferencia] Ejecución de un modelo entrenado sobre un audio de
    entrada para obtener una predicción o un resultado de rechazo según
    el predictor utilizado.

    \item[Ingesta] Incorporación de grabaciones al sistema mediante las
    operaciones de carga, importación o sincronización habilitadas.

    \item[Instantánea de entrenamiento (\textit{snapshot})] Selección
    identificada de grabaciones y metadatos que sirve como fuente de un
    experimento. La identificación de la fuente permite rastrear qué
    datos utilizó un trabajo.

    \item[Investigador] Usuario que trabaja con los conjuntos de datos que
    tiene autorizados. Puede preparar experimentos, entrenar modelos,
    consultar resultados y revisar predicciones dentro de su ámbito de acceso.

    \item[Macro F1] Promedio aritmético del F1 de las clases evaluadas.
    El F1 de cada clase combina su precisión y exhaustividad mediante la
    media armónica. El promedio asigna el mismo peso a cada clase.

    \item[Metadatos acústicos] Información que describe una grabación,
    como su duración, frecuencia de muestreo, origen, clase y ubicación,
    según los datos registrados.

    \item[Modelo activo] Modelo designado explícitamente como
    predeterminado en el ámbito correspondiente. Su utilización requiere
    que el predictor pueda acceder a los pesos necesarios.

    \item[Nodo de procesamiento local] Estación o servidor que ejecuta
    el procesamiento acústico y las operaciones de entrenamiento e
    inferencia de F.A.M.A.

    \item[Particiones] Subconjuntos que se reservan para entrenar, validar
    y probar el modelo. El entrenamiento ajusta los pesos, la validación
    orienta la selección del modelo y la prueba permite estimar su
    desempeño sobre ejemplos reservados.

    \item[Precisión (\textit{precision})] Proporción de predicciones de
    una clase que coinciden con la etiqueta de referencia de esa clase.

    \item[Preprocesamiento] Transformaciones que preparan la grabación
    para el modelo. Pueden incluir remuestreo, selección de ventanas,
    extracción de características y normalización.

    \item[Retroalimentación (\textit{feedback})] Confirmación o corrección
    que el usuario registra sobre una predicción. Un aporte incorporado
    puede incluirse en una nueva preparación para un entrenamiento posterior.

    \item[Trazabilidad] Capacidad de identificar la relación entre los
    datos y particiones utilizados, el modelo obtenido, su evaluación y
    las predicciones o correcciones asociadas.
\end{description}

\section{Abreviaciones}

\begin{description}
    \item[CP] Caso de prueba.
    \item[CPU] Unidad central de procesamiento.
    \item[CU] Caso de uso.
    \item[F.A.M.A.] Flujo de clasificación de Audio y Machine-learning Avanzado.
    \item[GCP] \textit{Google Cloud Platform}, plataforma de servicios en la nube.
    \item[GCS] \textit{Google Cloud Storage}, servicio de almacenamiento de objetos.
    \item[GPU] Unidad de procesamiento gráfico.
    \item[HPSS] Separación armónica y percusiva
    (\textit{Harmonic-Percussive Source Separation}).
    \item[IAM] Gestión de identidades y accesos
    (\textit{Identity and Access Management}).
    \item[MFCC] Coeficientes cepstrales en escala Mel
    (\textit{Mel-Frequency Cepstral Coefficients}).
    \item[MLOps] \textit{Machine Learning Operations}. Prácticas para
    organizar y operar el ciclo de vida de los modelos de aprendizaje
    automático.
    \item[MVP] Producto mínimo viable (\textit{Minimum Viable Product}).
    \item[RF] Requerimiento funcional.
    \item[RMS] Raíz cuadrática media (\textit{Root Mean Square}).
    \item[RNF] Requerimiento no funcional.
    \item[STFT] Transformada de Fourier de tiempo corto
    (\textit{Short-Time Fourier Transform}).
    \item[TDD] Desarrollo guiado por pruebas
    (\textit{Test-Driven Development}).
    \item[VAN] Valor actual neto.
    \item[VRAM] Memoria de video utilizada por la GPU.
\end{description}
